\chapter{Time Value of Money, Discounted Cash Flows, and Valuation Multiples}
\label{ch:valuation-multiples}
\chapterline

Part~I built a conservation-law model of accounting: double-entry ledgers record every economic event without losing a dollar, and the three financial statements summarize what a business owns, what it owes, and what it earned over past reporting periods. Accounting looks backward with precision. Valuation asks the forward-looking question: given an operating machine that will generate uncertain cash flows over the next ten, twenty, or fifty years, what is the maximum price a rational buyer should pay for it today?

Every serious answer to that question -- whether a multi-tab discounted cash flow spreadsheet, a venture partner's back-of-the-envelope revenue multiple, or a value investor's price-to-earnings ratio -- rests on a single mathematical conversion between cash today and cash received at a future date. We will derive that conversion from scratch, build discounted cash flow valuation on top of it, distinguish equity value from enterprise value, and prove algebraically when a high valuation multiple is rational and when a low multiple is a trap.

\begin{historybox}
The mathematics of present value grew directly out of commercial trade and sovereign borrowing. In 1202, Leonardo of Pisa (known later as Fibonacci) published \emph{Liber Abaci}, introducing Hindu-Arabic numerals to European merchants. Alongside currency conversion and alloy problems, Fibonacci posed concrete contract comparisons: which payment schedule is worth more to a merchant when florins arrive in staggered installments and capital can earn compound return in trade? Solving those comparisons required pulling future cash flows back to a single reference date.

Four and a half centuries later, in 1671, Johan de Witt -- mathematician and Grand Pensionary of Holland -- wrote \emph{Waerdye van Lyf-Renten Naer Proportie van Los-Renten} (\emph{Value of Life Annuities in Proportion to Redeemable Annuities}) to finance the Dutch Republic before the Franco-Dutch War. Combining mortality tables from municipal registers with compound discounting at $4\%$, de Witt proved that the state was selling life annuities far below their actuarial present value. In the twentieth century, Irving Fisher's \emph{The Theory of Interest} (1930) formalized asset prices as the discounted sum of future income streams, and Benjamin Graham and David Dodd's \emph{Security Analysis} (1934) -- written in the wreckage of the 1929 market crash -- separated a company's intrinsic value, anchored in cash-generating power and balance-sheet claims, from the daily quotations of an emotional market.
\end{historybox}

\section{The Time Value of Money and Net Present Value}

Adding \$100 received today to \$100 received five years from now is a dimensional error, like adding meters to yards without a conversion factor. Cash dated $t = 0$ and cash dated $t = 5$ are different economic goods for three independent reasons:

\begin{enumerate}
  \item \textbf{Real risk-free opportunity cost.} Even in a world with zero inflation and zero default risk, \$100 in hand today can be deployed into productive equipment or risk-free sovereign Treasury bills to become more than \$100 next year. Waiting five years forfeits five years of productive compounding.
  \item \textbf{Inflation.} When the general price level rises at annual rate $\pi$, a dollar delivered in year $t$ buys fewer kilowatt-hours, fewer hours of engineering labor, and fewer loaves of bread than a dollar today. If you require a real purchasing-power return $r_{\text{real}}$, the nominal rate $r_{\text{nom}}$ must satisfy Fisher's exact compounding relation,
  \[
    1 + r_{\text{nom}} = (1 + r_{\text{real}})(1 + \pi) = 1 + r_{\text{real}} + \pi + r_{\text{real}}\pi.
  \]
  \item \textbf{Default and equity risk premium.} A promised corporate coupon might default, and a projected software cash flow might shrink if a competitor ships a better product. Because human investors prefer predictable consumption over volatile consumption -- especially against losses that arrive during recessions when jobs and other assets are simultaneously impaired -- risky cash flows require an extra expected return $r_{\text{risk}} > 0$ above the risk-free government rate $r_f$.
\end{enumerate}

The combined annual rate $r = r_f + r_{\text{risk}}$ is the \textbf{discount rate} (or \textbf{opportunity cost of capital}): the expected annual return an investor gives up by putting capital into this asset instead of publicly traded securities of equivalent risk.

\subsection{Discrete compounding, continuous compounding, and discount factors}

Suppose you invest principal $P_0$ at annual rate $r$ compounded once per year. After one year you hold $P_0(1+r)$; after two years, $P_0(1+r)^2$; after $t$ years,
\[
  FV_t = P_0(1 + r)^t.
\]
If interest is compounded $m$ times per year (for example, $m = 4$ for quarterly compounding or $m = 365$ for daily compounding), each sub-period applies rate $r/m$ across $mt$ total periods:
\[
  FV_t^{(m)} = P_0\left(1 + \frac{r}{m}\right)^{mt} = P_0\left[\left(1 + \frac{r}{m}\right)^m\right]^t.
\]
Letting the compounding frequency grow without bound ($m \to \infty$) invokes the standard exponential limit $\lim_{m \to \infty}(1 + r/m)^m = e^r$, yielding \textbf{continuous compounding}:
\[
  FV_t^{(\infty)} = \lim_{m \to \infty} P_0\left(1 + \frac{r}{m}\right)^{mt} = P_0\,e^{rt}.
\]
At $r = 8\%$ over $t = 1$ year, \$1,000 grows to \$1,080.00 under annual compounding, \$1,082.43 under quarterly compounding ($m=4$), and \$1,083.29 under continuous compounding ($e^{0.08} \approx 1.083287$). Corporate valuation and accounting filings almost universally use discrete annual or quarterly compounding; quantitative derivatives pricing and continuous-time FX models use $e^{rt}$ because derivatives and integrals of $e^{rt}$ stay simple.

Valuation runs compounding in reverse. Ask: how much principal $PV$ must be invested today at annual rate $r$ so that it grows to match a future cash payment $C_t$ in year $t$? Solving $PV \cdot (1+r)^t = C_t$ gives
\[
  PV = \frac{C_t}{(1+r)^t} = C_t \cdot d_t, \qquad \text{where } d_t = (1+r)^{-t} \text{ (or } e^{-rt} \text{ continuously).}
\]
The multiplier $d_t = (1+r)^{-t}$ is the \textbf{discount factor} for horizon $t$. Because $r > 0$, we have $0 < d_t < 1$, and $d_t$ decays exponentially with time. At $r = 10\%$, a dollar promised in $5$ years is worth $1.10^{-5} \approx 0.621$ dollars today; in $20$ years, $0.149$ dollars; in $50$ years, less than one cent ($0.0085$).

\subsection{Net Present Value}

Once every future dollar is multiplied by its date-specific discount factor $d_t$, all terms share the same unit -- present-day dollars -- and can be summed directly.

\begin{definition}[Net Present Value]
Given a stream of net cash flows $C_0, C_1, \dots, C_n$ occurring at dates $t = 0, 1, \dots, n$ and an opportunity cost of capital $r > -1$, the \textbf{Net Present Value} ($\NPV$) is
\[
  \NPV(r) = \sum_{t=0}^{n} \frac{C_t}{(1+r)^t} = C_0 + \frac{C_1}{1+r} + \frac{C_2}{(1+r)^2} + \dots + \frac{C_n}{(1+r)^n}.
\]
\end{definition}

In a typical engineering or capital-budgeting decision, $C_0 < 0$ is an immediate cash outlay (building a factory line, acquiring a competitor, reserving compute hardware) and $C_1, \dots, C_n$ are the incremental net cash inflows generated by that decision.

\begin{example}[Evaluating an automation outlay]\label{ex:ch03-automation}
A hardware manufacturer can spend \$1,000,000 today ($C_0 = -1.00$ in USD millions) on automated optical inspection equipment that reduces scrap and warranty returns by \$350,000 per year for four years ($C_1 = C_2 = C_3 = C_4 = +0.35$ in USD millions), after which the line is retired with zero salvage value. For projects of this operating risk, the company's cost of capital is $r = 10\%$. Discounting each year's cash savings back to $t = 0$:
\begin{align*}
  \NPV(0.10) &= -1.00 + \frac{0.35}{1.10} + \frac{0.35}{1.10^2} + \frac{0.35}{1.10^3} + \frac{0.35}{1.10^4} \\
  &= -1.00 + 0.3182 + 0.2893 + 0.2630 + 0.2391 = +0.1095 \text{ (USD millions)}.
\end{align*}
Undertaking the project creates \$109,453 of net present wealth above what the \$1,000,000 could earn in equivalent-risk investments at $10\%$. Notice that undiscounted cash inflows summed to \$1,400,000 (an apparent \$400,000 profit), yet over three-quarters of that raw surplus merely compensated capital for waiting four years at $10\%$. Had the cost of capital been $16\%$, the exact same cash flows would yield $\NPV(0.16) = -1.00 + 0.9794 = -0.0206$ (a loss of \$20,636 in present value), and the rational choice would be to reject the project.
\end{example}

\begin{intuitionbox}
Why does $\NPV > 0$ serve as the clean decision rule for capital allocation? Because $\NPV$ obeys \textbf{value additivity}: $\NPV(A + B) = \NPV(A) + \NPV(B)$. Accepting a project with $\NPV =$ +\$109,453 is economically equivalent to receiving a \$109,453 cash check today while investing your \$1,000,000 at the fair market rate $r$. Any competing metric that violates value additivity can trick a decision-maker into rejecting wealth-increasing combinations or accepting wealth-destroying ones.
\end{intuitionbox}

\section{Internal Rate of Return and Its Mathematical Traps}

Executives and private-equity investors often prefer quoting a percentage rate of return rather than a dollar present value. Instead of plugging in a hurdle rate $r$ and computing $\NPV(r)$, they ask: \emph{at what break-even discount rate would this project's Net Present Value equal zero?}

\begin{definition}[Internal Rate of Return]
The \textbf{Internal Rate of Return} ($\IRR$) of a cash-flow stream $C_0, C_1, \dots, C_n$ is any rate $\IRR > -1$ satisfying
\[
  \NPV(\IRR) = \sum_{t=0}^{n} \frac{C_t}{(1+\IRR)^t} = 0.
\]
\end{definition}

Algebraically, substitute $x = \frac{1}{1+\IRR}$. The defining condition becomes a degree-$n$ polynomial equation in $x$:
\[
  P(x) = C_0 + C_1 x + C_2 x^2 + \dots + C_n x^n = 0.
\]
When does $\IRR$ behave cleanly? Suppose a project is \textbf{conventional}: cash flows start strictly negative ($C_0 < 0$) and are non-negative thereafter ($C_t \ge 0$ for all $t \ge 1$, with at least one $C_t > 0$). Differentiating $\NPV(r)$ term by term on $r \in (-1, \infty)$ gives
\[
  \dv{\NPV}{r} = -\sum_{t=1}^{n} \frac{t\,C_t}{(1+r)^{t+1}} < 0.
\]
Because $\NPV(r)$ is continuous and strictly decreasing on $(-1, \infty)$, diverges to $+\infty$ as $r \to -1^+$, and approaches $C_0 < 0$ as $r \to \infty$, the Intermediate Value Theorem guarantees a \textbf{unique} root $\IRR \in (-1, \infty)$. (Equivalently, Descartes' Rule of Signs confirms a single positive root $x > 0$ because the sequence $C_0, C_1, \dots, C_n$ changes sign exactly once.) For our optical inspection machine in Example~\ref{ex:ch03-automation}, solving $-1.00 + \sum_{t=1}^4 0.35(1+\IRR)^{-t} = 0$ numerically yields a unique $\IRR = 14.96\%$. Since $\NPV(r)$ slopes strictly downward, $\IRR > r$ holds if and only if $\NPV(r) > 0$.

Outside a single conventional project, relying on $\IRR$ without checking $\NPV$ walks straight into two classic mathematical failures.

\begin{trapbox}
\textbf{Trap 1: Non-unique (or non-existent) roots when cash flows change sign twice.}
Consider an energy extraction, hardware-leasing, or infrastructure concession project requiring \$100 million upfront ($C_0 = -100$), generating \$230 million of cash inflow in year 1 ($C_1 = +230$), and requiring \$132 million of mandatory environmental remediation or decommissioning in year 2 ($C_2 = -132$). Setting $\NPV(\IRR) = 0$ and multiplying through by $(1+\IRR)^2$ gives
\[
  -100(1+\IRR)^2 + 230(1+\IRR) - 132 = 0 \iff \bigl((1+\IRR) - 1.10\bigr)\bigl((1+\IRR) - 1.20\bigr) = 0.
\]
This single project has \textbf{two exact Internal Rates of Return}: $\IRR_1 = 10\%$ and $\IRR_2 = 20\%$. Worse, the parabola opens downward: $\NPV(r)$ is strictly positive \emph{between} $10\%$ and $20\%$ (for instance, at $r = 15\%$, $\NPV = -100 + 200 - 99.81 =$ +\$0.19 million) and strictly negative for both $r < 10\%$ and $r > 20\%$ (at $r = 5\%$, $\NPV =$ -\$0.68 million because the low discount rate barely shrinks the year-2 remediation liability). Asking ``is $\IRR$ above our $12\%$ hurdle rate?'' has no coherent meaning when $\IRR$ is simultaneously $10\%$ and $20\%$. Evaluating $\NPV(0.12) =$ +\$0.19 million resolves the decision immediately without ambiguity.

\medskip
\textbf{Trap 2: Scale blindness across mutually exclusive choices.}
Suppose you have rights to build on a single parcel of land and must choose one of two mutually exclusive projects at a $10\%$ cost of capital:
\begin{itemize}
  \item \textbf{Project A (kiosk):} Spend \$1,000 today, receive \$2,000 in one year. Here $\IRR_A = 100\%$, and $\NPV_A(0.10) = -1{,}000 + \frac{2{,}000}{1.10} =$ +\$818.18.
  \item \textbf{Project B (warehouse):} Spend \$1,000,000 today, receive \$1,350,000 in one year. Here $\IRR_B = 35\%$, and $\NPV_B(0.10) = -1{,}000{,}000 + \frac{1{,}350{,}000}{1.10} =$ +\$227,272.73.
\end{itemize}
Ranking by $\IRR$ picks the $100\%$ return over the $35\%$ return and forfeits \$226,454.55 of net present wealth. Percentage rates carry no information about how much capital can actually be put to work at that rate (nor how long it compounds). When comparing mutually exclusive projects, maximize dollar $\NPV$, or verify that the incremental outlay $B - A$ ($C_0 =$ -\$999,000, $C_1 =$ +\$1,348,000, incremental $\IRR = 34.93\% > 10\%$) earns well above the cost of capital.
\end{trapbox}

\section{Perpetuities, Gordon Growth, and Two-Stage DCF Valuation}

A factory machine wears out in four years, whereas a profitable corporation can hire new engineers, launch replacement products, and generate \textbf{Free Cash Flow} ($\FCF$ -- the operating cash left over after maintaining and expanding the asset base, defined formally in Part~I and dissected in Chapter~\ref{ch:filings-and-profitability}) for decades without a fixed shutdown date. How do we discount an infinite sequence of growing cash flows into a finite present price?

\subsection{Exact derivation of the growing perpetuity}

Suppose a business generates free cash flow $\FCF_0$ in the base year just ended, and its cash flows grow at constant annual rate $g$ forever, so that year $t$ delivers $\FCF_t = \FCF_0(1+g)^t = \FCF_1(1+g)^{t-1}$, while investors discount those cash flows at rate $r > 0$.

\begin{proposition}[Gordon Growth / Growing Perpetuity Formula]\label{prop:ch03-gordon}
Let $\FCF_1 = \FCF_0(1+g)$ be next year's expected cash flow, growing at constant rate $g > -1$ in perpetuity and discounted at rate $r > -1$. The infinite present-value sum converges to a finite value if and only if $r > g$, in which case
\begin{equation}
  V_0 = \sum_{t=1}^{\infty} \frac{\FCF_0(1+g)^t}{(1+r)^t} = \frac{\FCF_0(1+g)}{r - g} = \frac{\FCF_1}{r - g}.
  \label{eq:gordon-growth}
\end{equation}
\end{proposition}

\begin{proof}
Define the growth-to-discount ratio $q = \frac{1+g}{1+r} > 0$, and write the partial sum of the first $N$ years as
\[
  S_N = \sum_{t=1}^{N} \FCF_0 \left(\frac{1+g}{1+r}\right)^t = \FCF_0 \sum_{t=1}^{N} q^t.
\]
Multiply $S_N$ by $q$ and subtract from $S_N$ so all interior terms telescope:
\[
  S_N - q S_N = \FCF_0\bigl(q + q^2 + \dots + q^N\bigr) - \FCF_0\bigl(q^2 + \dots + q^N + q^{N+1}\bigr) = \FCF_0\bigl(q - q^{N+1}\bigr).
\]
If $q \ge 1$ (meaning $g \ge r$), then every term $\FCF_0 q^t \ge \FCF_0 > 0$, so $S_N \ge N \cdot \FCF_0 \to \infty$ as $N \to \infty$ (assuming $\FCF_0 > 0$). If $0 < q < 1$ (which holds if and only if $1 + g < 1 + r$, that is, $r > g$), we divide by $1 - q \neq 0$:
\[
  S_N = \FCF_0\,\frac{q(1 - q^N)}{1 - q}.
\]
Taking the limit as $N \to \infty$, $0 < q < 1$ implies $\lim_{N \to \infty} q^N = 0$. Meanwhile, simplifying $1 - q$ over a common denominator gives
\[
  1 - q = 1 - \frac{1+g}{1+r} = \frac{(1+r) - (1+g)}{1+r} = \frac{r - g}{1+r}.
\]
Substituting $q = \frac{1+g}{1+r}$ and $1 - q = \frac{r - g}{1+r}$ into the limit yields
\[
  V_0 = \lim_{N \to \infty} S_N = \FCF_0 \cdot \frac{\dfrac{1+g}{1+r}}{\dfrac{r - g}{1+r}} = \frac{\FCF_0(1+g)}{r - g} = \frac{\FCF_1}{r - g}. \qedhere
\]
\end{proof}

Notice two structural boundary cases:
\begin{itemize}
  \item Setting $g = 0$ recovers the \textbf{level perpetuity} $V_0 = \frac{\FCF_1}{r}$, first used to price British \emph{consols} and Dutch perpetual bonds: a security paying \$40 per year forever at $r = 5\%$ is worth $40 / 0.05 =$ \$800.
  \item Why must $g < r$ hold economically in perpetuity? No single company can compound faster than nominal global GDP forever -- otherwise that single firm would eventually exceed $100\%$ of world economic output. Long-run risk-free rates track nominal GDP growth, and risky equities carry a positive risk premium above the risk-free rate, forcing steady-state $g_\infty$ strictly below $r$.
\end{itemize}

\subsection{Two-stage Discounted Cash Flow (DCF) and terminal-value sensitivity}

Young, rapidly expanding companies routinely grow cash flows at $20\%$ or $35\%$ per year for five to ten years -- well above $r = 10\%$ -- before market saturation and competition pull growth down toward economy-wide rates. Plugging $g = 25\%$ and $r = 10\%$ into \eqref{eq:gordon-growth} produces a nonsensical negative number because the series diverges. Instead, analysts split time into two stages:

\begin{enumerate}
  \item \textbf{Explicit forecast stage ($t = 1, \dots, T$, typically $5$--$10$ years):} Model revenue, margins, working capital, and capital expenditures year by year while growth and profitability are still changing, obtaining $\FCF_1, \dots, \FCF_T$.
  \item \textbf{Steady-state terminal stage ($t > T$):} Once the company reaches maturity at year $T$ with sustainable long-run growth $g_\infty < r$, apply \eqref{eq:gordon-growth} standing at date $T$ to value all cash flows from year $T+1$ onward:
  \[
    V_T = \sum_{k=1}^{\infty} \frac{\FCF_T(1 + g_\infty)^k}{(1+r)^k} = \frac{\FCF_{T+1}}{r - g_\infty} = \frac{\FCF_T(1 + g_\infty)}{r - g_\infty}.
  \]
  Because $V_T$ is measured in year-$T$ dollars, we multiply it by $(1+r)^{-T}$ to bring it back to $t = 0$.
\end{enumerate}

Combining both stages gives the canonical \textbf{two-stage DCF equation}:
\begin{equation}
  V_0 = \underbrace{\sum_{t=1}^{T} \frac{\FCF_t}{(1+r)^t}}_{\text{PV of explicit forecast}} + \underbrace{\frac{1}{(1+r)^T} \cdot \frac{\FCF_T(1 + g_\infty)}{r - g_\infty}}_{\text{PV of terminal value}}.
  \label{eq:two-stage-dcf}
\end{equation}

\begin{figure}[htbp]
\centering
\begin{tikzpicture}
  % Panel (a): Present value decay of USD 100 across discount rates
  \begin{scope}[shift={(0,0)}]
    \begin{axis}[
      width=5.7cm, height=5.1cm,
      title={Panel (a): PV of \$100 in year $t$},
      title style={font=\footnotesize\bfseries, text=sourceblue},
      xlabel={horizon $t$ (years)},
      ylabel={present value (USD)},
      label style={font=\footnotesize},
      tick label style={font=\scriptsize},
      xmin=0, xmax=30, ymin=0, ymax=106,
      axis lines=left, samples=61, domain=0:30,
      xtick={0,10,20,30}, ytick={0,25,50,75,100},
      enlargelimits=false, clip=false
    ]
      \addplot[targetgreen, very thick] {100*(1.04)^(-x)};
      \addplot[curveblue, very thick, dashed] {100*(1.08)^(-x)};
      \addplot[trapred, very thick, densely dotted] {100*(1.12)^(-x)};
      \node[targetgreen, font=\scriptsize, right] at (axis cs:18,54) {$r=4\%$};
      \node[curveblue, font=\scriptsize, right] at (axis cs:14,37) {$r=8\%$};
      \node[trapred, font=\scriptsize, right] at (axis cs:9,21) {$r=12\%$};
    \end{axis}
  \end{scope}
  % Panel (b): Two-stage DCF cash flows (Nominal vs PV)
  \begin{scope}[shift={(6.4,0)}]
    \begin{axis}[
      width=5.8cm, height=5.1cm,
      title={Panel (b): Nominal vs.\ PV $\FCF_t$ ($r=10\%$)},
      title style={font=\footnotesize\bfseries, text=targetgreen},
      xlabel={year $t$ (Stage 1: $1$--$5$, Stage 2: $6+$)},
      ylabel={$\FCF$ (USD millions)},
      label style={font=\footnotesize},
      tick label style={font=\scriptsize},
      xmin=0.4, xmax=12.6, ymin=0, ymax=150,
      axis lines=left,
      xtick={1,3,5,7,9,11}, ytick={0,50,100,150},
      enlargelimits=false, clip=false
    ]
      % Nominal FCF bars (light background)
      \addplot[ybar, bar width=6pt, draw=quotientorange!80!black, fill=quotientorange!25] coordinates {
        (1,59.0) (2,69.6) (3,82.2) (4,96.9) (5,114.4)
        (6,117.8) (7,121.4) (8,125.0) (9,128.7) (10,132.6) (11,136.6) (12,140.7)
      };
      % Present Value FCF bars (darker foreground overlay)
      \addplot[ybar, bar width=3.5pt, draw=sourceblue!90!black, fill=sourceblue!70] coordinates {
        (1,53.6) (2,57.5) (3,61.7) (4,66.2) (5,71.0)
        (6,66.5) (7,62.3) (8,58.3) (9,54.6) (10,51.1) (11,47.9) (12,44.8)
      };
      \draw[gray!75!black, dashed, thick] (axis cs:5.5,0) -- (axis cs:5.5,145);
      \node[quotientorange!85!black, font=\scriptsize, right] at (axis cs:0.7,134) {Nominal $\FCF_t$};
      \node[sourceblue!90!black, font=\scriptsize, right] at (axis cs:0.7,117) {PV $\FCF_t/(1+r)^t$};
      \node[gray!60!black, font=\tiny, right] at (axis cs:5.65,140) {$g_\infty=3\%$};
    \end{axis}
  \end{scope}
\end{tikzpicture}
\caption{Panel (a) shows exponential decay of \$100 promised at horizon $t$ under $r = 4\%, 8\%, 12\%$. Panel (b) plots a two-stage DCF starting from $\FCF_0 =$ \$50M with $18\%$ growth in years $1$--$5$ and $3\%$ steady-state growth thereafter at $r = 10\%$. While nominal cash flows (tan) climb forever, discounted bars (blue) rise only while $g > r$ ($t \le 5$) and decay geometrically once $g_\infty < r$ ($t \ge 6$), ensuring a finite sum.}
\label{fig:ch03-pv-and-dcf}
\end{figure}

Figure~\ref{fig:ch03-pv-and-dcf}(b) traces a concrete business starting at $\FCF_0 =$ \$50.0 million that compounds at $18\%$ for $T = 5$ years (reaching $\FCF_5 =$ \$114.39 million) and $g_\infty = 3\%$ thereafter at $r = 10\%$. Summing the five discounted blue bars to the left of the dashed line gives \$310.14 million. At $t = 5$, $\FCF_6 = 114.39 \times 1.03 =$ \$117.82 million, so $V_5 = 117.82 / (0.10 - 0.03) =$ \$1,683.14 million, whose present value at $t = 0$ is $1{,}683.14 / 1.10^5 =$ \$1,045.10 million. Total intrinsic value is $V_0 = 310.14 + 1{,}045.10 =$ \$1,355.24 million -- with $77.1\%$ of the entire valuation sitting inside the terminal value.

Because $r - g$ sits in the denominator of \eqref{eq:gordon-growth}, valuation is acutely sensitive to small shifts in discount rates or perpetual growth assumptions. Differentiating $V_0 = \frac{\FCF_1}{r - g}$ with respect to $r$ gives
\[
  \pdv{V_0}{r} = -\frac{\FCF_1}{(r - g)^2} = -\frac{V_0}{r - g}.
\]
When $r - g = 0.04$ (for example, $r = 8\%$ and $g = 4\%$), a $1$ percentage-point increase in interest rates ($\Delta r = +0.01$) cuts valuation by roughly $1 / 0.05$ vs.\ $1 / 0.04$, an exact $20\%$ drop -- explaining why long-duration growth stocks sell off sharply whenever central banks raise long-term yields. Table~\ref{tab:ch03-gordon-sensitivity} tabulates $V_0$ for $\FCF_1 =$ \$100 million across realistic $(r, g)$ pairs: moving from $(r = 10\%, g = 2\%)$ to $(r = 8\%, g = 4\%)$ doubles the company's value from \$1,250 million to \$2,500 million without changing next year's cash flow by a single cent.

\begin{table}[htbp]
\centering
\caption{Present value $V_0 = \FCF_1 / (r - g)$ (in USD millions) and implied cash-flow multiple ($V_0 / \FCF_1$) for $\FCF_1 =$ \$100 million.}
\label{tab:ch03-gordon-sensitivity}
\begin{tabular}{lccc}
\toprule
Discount rate $r$ & $g = 2.0\%$ & $g = 3.0\%$ & $g = 4.0\%$ \\
\midrule
$r = 8.0\%$  & \$1,667M ($16.7\times$) & \$2,000M ($20.0\times$) & \$2,500M ($25.0\times$) \\
$r = 9.0\%$  & \$1,429M ($14.3\times$) & \$1,667M ($16.7\times$) & \$2,000M ($20.0\times$) \\
$r = 10.0\%$ & \$1,250M ($12.5\times$) & \$1,429M ($14.3\times$) & \$1,667M ($16.7\times$) \\
\bottomrule
\end{tabular}
\end{table}

\section{Market Capitalization versus Enterprise Value}

When financial news headlines state that a public company is ``worth \$500 million,'' they are almost always quoting its \textbf{Market Capitalization} (or \textbf{Equity Value}):
\[
  \text{Market Cap} = \text{Share Price} \times \text{Diluted Common Shares Outstanding}.
\]
Market Cap is the price tag of the common equity slice alone. Yet as we saw on the Balance Sheet in Chapter~2, common shareholders stand at the back of the line behind lenders and preferred stockholders, while also owning whatever excess cash sits inside the corporate bank accounts. To price the underlying operating machine -- factories, servers, intellectual property, customer contracts -- independent of how it was financed, we need \textbf{Enterprise Value} ($\EV$).

Derive $\EV$ by asking a concrete takeover question: \emph{what would it actually cost you out of pocket today to buy $100\%$ of this operating business free and clear of all financial claims, just like buying a house free and clear of its mortgage?}

\begin{enumerate}
  \item First, you write a check equal to \textbf{Market Cap} to buy out $100\%$ of common stockholders.
  \item Second, to own the operating assets unencumbered, you must pay off (or assume legal liability for) all \textbf{Total Debt} (short-term debt, current maturities, long-term bonds, bank term loans, and finance leases) as well as any \textbf{Preferred Stock} (senior equity-like claims with fixed liquidation value; if the company consolidates a partially owned subsidiary, you also add \textbf{Non-Controlling / Minority Interest}).
  \item Third, the moment the acquisition closes, you gain full control of the company's \textbf{Cash and Cash Equivalents} (unrestricted bank cash and short-term Treasury bills sitting on the balance sheet), which you can wire straight back to yourself to reimburse part of your purchase check.
\end{enumerate}

Summing the checks written minus the cash recovered inside the vault yields the fundamental identity:
\begin{equation}
  \EV = \text{Market Cap} + \text{Total Debt} + \text{Preferred Stock} - \text{Cash \& Equivalents}.
  \label{eq:enterprise-value}
\end{equation}

Compare two concrete takeover candidates side by side to see why looking at Market Cap alone misleads beginners:
\begin{itemize}
  \item \textbf{Company A (cash-rich software firm):} Market Cap is \$500 million, Total Debt is \$40 million, Preferred Stock is \$0, and Cash \& Equivalents is \$150 million. Buying $100\%$ of the equity for \$500 million hands you \$150 million of balance-sheet cash and \$40 million of debt obligations. Your true net takeover price for the software business is
  \[
    \EV_A = 500 + 40 + 0 - 150 = 390 \text{ (USD millions)}.
  \]
  \item \textbf{Company B (debt-laden industrial firm):} Market Cap is just \$100 million -- five times ``smaller'' on a stock-screener headline -- with \$0 excess cash, \$0 preferred stock, and \$400 million of senior debt. Acquiring the operating business free and clear actually costs
  \[
    \EV_B = 100 + 400 + 0 - 0 = 500 \text{ (USD millions)}.
  \]
\end{itemize}
The ``cheap'' \$100 million stock is a \$500 million enterprise purchase, more expensive than the \$500 million software company whose operating business costs \$390 million.

\begin{formalbox}
\textbf{Consistency law for valuation ratios (matching numerator and denominator):}
\begin{itemize}
  \item Any financial metric measured \emph{before} interest expense -- \textbf{Revenue}, \textbf{Gross Profit}, \textbf{EBITDA}, \textbf{EBIT}, or \textbf{Unlevered $\FCF$} -- flows to \emph{all} providers of capital (lenders and shareholders alike) and must be paired with \textbf{Enterprise Value ($\EV$)}.
  \item Any financial metric measured \emph{after} interest expense and preferred dividends -- \textbf{Net Income}, \textbf{Earnings Per Share (EPS)}, or \textbf{Levered $\FCF$ to Equity} -- flows exclusively to common shareholders and must be paired with \textbf{Market Cap} (or \textbf{Share Price}).
\end{itemize}
Forming ``Price-to-Sales'' ($\text{Market Cap} / \text{Revenue}$) on a heavily indebted retailer makes the stock look artificially cheap simply because debt replaced equity in the numerator while $100\%$ of store revenue remained in the denominator. Always use $\EV / \text{Revenue}$.
\end{formalbox}

\section{Valuation Multiples Derived From First Principles}

In practice, traders, board members, and acquirers constantly quote ratios like ``$22\times$ forward earnings'' or ``$11\times$ $\EV/\text{EBITDA}$.'' Beginners often treat multiples as folk rules of thumb disconnected from discounted cash flows. In fact, every sensible multiple is a compressed algebraic rearrangement of a DCF equation -- and writing out that algebra exposes the hidden growth, return-on-capital, and risk assumptions baked into every multiple.

\subsection{The P/E ratio, TTM versus Forward P/E, and Earnings Yield}

The \textbf{Price-to-Earnings (P/E) ratio} divides the market value of common equity by the net profit attributable to common shareholders:
\[
  \text{P/E} = \frac{\text{Market Cap}}{\text{Net Income}} = \frac{\text{Share Price } P_0}{\text{Diluted EPS}}.
\]
Two conventions appear across financial terminals and analyst reports:
\begin{itemize}
  \item \textbf{Trailing Twelve Months (TTM) P/E} divides today's price $P_0$ by actual reported EPS summed over the four most recent completed quarters ($E_0$). Its denominator is audited or reported fact, though backward-looking when a business is growing or entering a downturn.
  \item \textbf{Forward (Next Twelve Months, NTM) P/E} divides today's price $P_0$ by projected EPS over the upcoming four quarters ($E_1$). Because asset prices discount \emph{future} cash flows, Forward P/E matches the timing of \eqref{eq:gordon-growth}, at the cost of relying on analyst or management forecasts.
\end{itemize}

Look closely at the units of P/E: Share Price has units of $\text{USD}$, and annual EPS has units of $\text{USD} / \text{year}$, so $\text{P/E}$ has units of \textbf{years} (how many years of current earnings equal today's stock price). Inverting the ratio gives the \textbf{Earnings Yield}:
\[
  \text{Earnings Yield} = \frac{E_1}{P_0} = \frac{1}{\text{Forward P/E}} \quad (\text{units: } \% \text{ per year}).
\]
Inverting P/E into an annual percentage yield allows direct comparison against government bond yields. Suppose the 10-year U.S. Treasury bond pays a risk-free coupon yield of $4.5\%$ ($P / \text{Coupon} = 22.2\times$). A mature utility or consumer company trading at $20\times$ forward earnings offers a $5.0\%$ initial earnings yield ($1/20 = 0.05$). Unlike a fixed Treasury coupon, corporate earnings can grow with inflation and reinvested capital -- or shrink during a recession. Comparing $E_1 / P_0$ against the sovereign bond yield immediately poses the central equity question: does the company's future earnings growth compensate for its business risk relative to locking in $4.5\%$ risk-free?

\subsection{Algebraic proof of what drives the P/E ratio}

What makes one company rationally deserve a $45\times$ forward P/E while another company is overpriced at $9\times$? Many investors casually answer ``earnings growth.'' Let us derive the exact formula for justified Forward P/E from first principles and watch why growth alone is only half the equation -- and sometimes contributes zero value.

Consider an unlevered business (or equity claim) expected to earn $E_1$ next year before new expansion investments. Each year, management retains and plows back a fraction $b \in [0, 1)$ of earnings into net new invested capital (capital expenditures above depreciation plus incremental working capital) and distributes the remaining fraction $1 - b$ to owners as free cash flow:
\[
  \FCF_1 = E_1(1 - b).
\]
Suppose each dollar of net new capital reinvested into the business earns a perpetual after-tax \textbf{Return on Invested Capital} ($\ROIC$). Reinvesting $b E_1$ dollars at date $1$ adds $\ROIC \cdot (b E_1)$ of permanent incremental annual earnings starting in year $2$. Therefore, annual earnings (and cash flows, for constant $b$ and $\ROIC$) compound at the exact rate
\begin{equation}
  g = \frac{\Delta E}{E_1} = \frac{\ROIC \cdot (b E_1)}{E_1} = b \cdot \ROIC.
  \label{eq:growth-plowback-roic}
\end{equation}

\begin{proposition}[Fundamental Value-Driver Equation for P/E]
For a firm reinvesting fraction $b \in [0, 1)$ of earnings at constant return on new invested capital $\ROIC$ with discount rate $r > b \cdot \ROIC$, the justified forward P/E ratio is
\begin{equation}
  \frac{P_0}{E_1} = \frac{1 - b}{r - b \cdot \ROIC} = \frac{1 - g / \ROIC}{r - g}.
  \label{eq:justified-pe}
\end{equation}
\end{proposition}

\begin{proof}
Substitute $\FCF_1 = E_1(1 - b)$ and $g = b \cdot \ROIC$ into Proposition~\ref{prop:ch03-gordon} ($P_0 = \frac{\FCF_1}{r - g}$):
\[
  P_0 = \frac{E_1(1 - b)}{r - b \cdot \ROIC}.
\]
Dividing both sides by $E_1 > 0$ gives $\frac{P_0}{E_1} = \frac{1 - b}{r - b \cdot \ROIC}$. For $\ROIC > 0$, substituting $b = g / \ROIC$ from \eqref{eq:growth-plowback-roic} yields the equivalent growth-explicit form $\frac{P_0}{E_1} = \frac{1 - g/\ROIC}{r - g}$.
\end{proof}

Equation~\eqref{eq:justified-pe} proves two facts that overturn naive multiple screening:

\medskip
\noindent\textbf{Fact 1: When $\ROIC = r$, faster growth creates zero value.}
Set $\ROIC = r$ directly inside \eqref{eq:justified-pe}:
\[
  \frac{P_0}{E_1} = \frac{1 - b}{r - b r} = \frac{1 - b}{r(1 - b)} = \frac{1}{r} \qquad \text{for every } b \in [0, 1).
\]
The plowback rate $b$ -- and therefore the growth rate $g = br$ -- cancels out identically. If $r = 10\%$ and a company earns $\ROIC = 10\%$ on new factories or acquisitions, its justified P/E is $1 / 0.10 = 10.0\times$ whether it pays out $100\%$ of earnings with $0\%$ growth ($b = 0$) or retains $75\%$ of earnings to grow at $7.5\%$ per year ($b = 0.75$). Why? Every \$1 withheld from shareholders today reduces current cash payout by \$1 and generates a perpetual stream of \$0.10 per year starting next year, whose present value at $r = 10\%$ is $0.10 / 0.10 =$ \$1.00 -- trading a dollar bill for four quarters. Worse, if $\ROIC < r$ (say $\ROIC = 6\%$ while $r = 10\%$), increasing $b$ to chase growth \emph{lowers} the justified P/E below $10.0\times$, because management is pouring shareholder dollars into projects worth \$0.60 on the dollar (Figure~\ref{fig:ch03-pe-roic}). Growth adds value if and only if $\ROIC > r$.

\medskip
\noindent\textbf{Fact 2: Why a $45\times$ P/E can be mathematically cheap while a $9\times$ P/E can be an overpriced value trap.}
Test both extremes numerically at the same $r = 10\%$ discount rate:
\begin{itemize}
  \item \textbf{High-ROIC compounder at $45\times$:} Suppose a capital-light software, payments network, or semiconductor IP firm earns $\ROIC = 30\%$ on incremental investments and has a long runway to reinvest $b = 28\%$ of earnings, compounding at $g = 0.28 \times 30\% = 8.4\%$. Equation~\eqref{eq:justified-pe} gives
  \[
    \frac{P_0}{E_1} = \frac{1 - 0.28}{0.10 - 0.084} = \frac{0.72}{0.016} = 45.0\times.
  \]
  Because $\ROIC = 30\%$ is triple the cost of capital, the company captures $8.4\%$ growth while still distributing $72\%$ of earnings as free cash flow. (In finite-horizon models where $\ROIC = 35\%$ allows $20\%$ annual compounding for $12$ years before fading to $3\%$, justified current P/E ratios readily exceed $40\times$--$50\times$.)
  \item \textbf{Two $9\times$ value traps:} First, consider a shrinking legacy business paying out all earnings ($b = 0$) whose revenue and net income erode by $5\%$ per year ($g = -0.05$) due to technological obsolescence. Its justified P/E is $\frac{1 - 0}{0.10 - (-0.05)} = \frac{1}{0.15} = 6.67\times$; paying $9.0\times$ overpays by $35\%$. Second, consider a capital-intensive manufacturer earning $\ROIC = 5\%$ that must plow back $b = 60\%$ of accounting profits into heavy machinery just to grow $g = 0.60 \times 5\% = 3.0\%$. Its justified P/E is $\frac{1 - 0.60}{0.10 - 0.03} = \frac{0.40}{0.07} = 5.71\times$; paying $9.0\times$ loses money because $60\%$ of reported ``earnings'' are permanently trapped in low-return metal and never reach shareholders.
\end{itemize}

\begin{figure}[htbp]
\centering
\begin{tikzpicture}
  \begin{axis}[
    width=10.6cm, height=6.0cm,
    title={Justified Forward P/E $\dfrac{P_0}{E_1} = \dfrac{1 - g/\ROIC}{r - g}$ vs.\ Growth Rate $g$ (at $r = 10\%$)},
    title style={font=\small\bfseries, text=sourceblue},
    xlabel={perpetual earnings growth rate $g$ (\%)},
    ylabel={justified forward P/E ($P_0 / E_1$)},
    label style={font=\footnotesize},
    tick label style={font=\footnotesize},
    xmin=0, xmax=8.0, ymin=0, ymax=38,
    axis lines=left, samples=100,
    xtick={0,1,2,3,4,5,6,7,8},
    ytick={0,10,20,30},
    grid=major, grid style={gray!20},
    enlargelimits=false, clip=false
  ]
    % ROIC = 30% > r (domain 0 to 8.0)
    \addplot[targetgreen, very thick, domain=0:8.0] {(1 - x/30) / (0.10 - x/100)};
    % ROIC = 10% = r (flat line at 10x)
    \addplot[curveblue, very thick, dashed, domain=0:8.0] {10};
    % ROIC = 6% < r (domain 0 to 6.0, hits 0 at g=6% where b=100%)
    \addplot[trapred, very thick, densely dotted, domain=0:6.0] {(1 - x/6) / (0.10 - x/100)};

    \node[targetgreen!85!black, font=\footnotesize\bfseries, left] at (axis cs:7.6,34.5) {$\ROIC = 30\% > r$ (value creation)};
    \node[curveblue, font=\footnotesize\bfseries, above right] at (axis cs:3.6,10.4) {$\ROIC = 10\% = r$ (flat at $1/r = 10\times$)};
    \node[trapred, font=\footnotesize\bfseries, below right] at (axis cs:1.8,5.8) {$\ROIC = 6\% < r$ (growth destroys value)};
    \filldraw[sourceblue] (axis cs:0,10) circle (2pt);
    \node[sourceblue, font=\scriptsize, right] at (axis cs:0.15,12.2) {$g=0 \implies 1/r = 10\times$};
  \end{axis}
\end{tikzpicture}
\caption{Justified Forward P/E as a function of earnings growth $g$ for fixed cost of capital $r = 10\%$. All three businesses start at $1/r = 10\times$ when $g = 0$. Growth lifts the multiple exclusively when $\ROIC > r$ (green), leaves the multiple completely invariant when $\ROIC = r$ (blue dashed), and drives the multiple toward zero when $\ROIC < r$ (red dotted, reaching $0\times$ at $g = 6\%$ where $b = 100\%$ of earnings are consumed by reinvestment).}
\label{fig:ch03-pe-roic}
\end{figure}

\subsection{Leverage-neutral multiples: EV/EBITDA, EV/FCF, and EV/Revenue}

Because P/E sits at the equity level, changing how a company finances itself distorts P/E even when the underlying factory or codebase is unchanged. Verify this on two identical operating businesses -- Firm~U (unlevered) and Firm~L (levered) -- each having Enterprise Value $\EV =$ \$1,000 million, $\text{EBITDA} =$ \$125 million, Depreciation \& Amortization $\text{D\&A} =$ \$25 million, operating profit $\text{EBIT} =$ \$100 million, and corporate tax rate $\tau = 25\%$:
\begin{itemize}
  \item \textbf{Firm U (\$0 debt, \$1,000M equity):} Interest expense is \$0, so pre-tax income is \$100M and Net Income is $100 \times (1 - 0.25) =$ \$75.0M. Its P/E ratio is $1{,}000 / 75.0 = 13.3\times$.
  \item \textbf{Firm L (\$500M debt at $6\%$ interest, \$500M equity):} Interest expense is $500 \times 0.06 =$ \$30M, so pre-tax income is $100 - 30 =$ \$70M and Net Income is $70 \times (1 - 0.25) =$ \$52.5M. Its P/E ratio is $500 / 52.5 = 9.5\times$.
\end{itemize}
A naive screener flags Firm~L ($9.5\times$ P/E) as $29\%$ cheaper than Firm~U ($13.3\times$ P/E), when in fact Firm~L owns the exact same operating machine at the exact same $\EV =$ \$1,000M -- its equity slice simply carries higher financial leverage and higher bankruptcy variance (causing its equity cost of capital $r_E$ to rise above Firm~U's).

Pairing Enterprise Value with pre-interest operating cash metrics eliminates this leverage distortion:
\begin{enumerate}
  \item \textbf{$\EV / \text{EBITDA}$:} Both Firm~U and Firm~L trade at the identical multiple $\EV / \text{EBITDA} = 1{,}000 / 125 = 8.0\times$ (and $\EV / \text{EBIT} = 10.0\times$). Because $\text{EBITDA} = \text{EBIT} + \text{D\&A}$ strips out interest, tax-rate differences, and historical depreciation schedules, it serves as the standard baseline multiple in corporate M\&A and leveraged credit covenants.
  \item \textbf{$\EV / \text{Unlevered }\FCF$:} EBITDA is not free cash flow whenever a business must pour real cash back into capital equipment or working capital. Defining unlevered free cash flow as $\FCF_U = \text{EBIT}(1 - \tau) + \text{D\&A} - \text{CapEx} - \Delta\text{NWC}$ produces $\EV / \FCF_U$, the purest single-period multiple of whole-firm cash generation.
  \item \textbf{$\EV / \text{Revenue}$:} Early high-growth cloud, marketplace, or AI companies routinely spend $40\%$--$80\%$ of revenue on R\&D and customer acquisition, rendering current EBITDA and Net Income zero or negative. Analysts quote $\EV / \text{Revenue}$ (for example, $10\times$ NTM revenue) as a bridge to future steady-state margins: if a software company trading at $10\times$ $\EV / \text{Revenue}$ has $80\%$ gross margins and will converge to a $25\%$ steady-state unlevered $\FCF$ margin at maturity, $10\times$ revenue corresponds to $10 / 0.25 = 40\times$ steady-state $\FCF$ (before accounting for intervening revenue compounding). Applying the same $10\times$ $\EV / \text{Revenue}$ multiple to a low-margin hardware reseller with $12\%$ gross margins and $3\%$ peak $\FCF$ margins implies an absurd $333\times$ steady-state $\FCF$. Never read an $\EV / \text{Revenue}$ multiple without multiplying through by gross margin and steady-state cash conversion.
\end{enumerate}

\begin{trapbox}
\textbf{Four classic multiple traps in live markets:}
\begin{enumerate}
  \item \textbf{Cyclical peak earnings.} Commodity producers -- steel mills, container shippers, energy refiners, and DRAM/NAND memory-chip fabs -- experience periodic supply shortages where selling prices double and operating margins surge $4\times$ above normal. At the top of the cycle, record TTM EPS drives the headline P/E down to $4\times$--$6\times$, luring beginners right before new fab capacity comes online, prices revert, and EPS collapses to zero (pushing trough P/E to $80\times$ or undefined precisely when the factories trade below replacement cost). Value cyclicals on mid-cycle normalized earnings or tangible book value, never peak TTM EPS.
  \item \textbf{One-off non-operating gains in GAAP Net Income.} Selling a headquarters building, divesting a division, marking up a minority equity stake, or releasing a deferred tax valuation allowance can double single-year GAAP Net Income and cut headline TTM P/E in half. One-time asset sales do not repeat in year $2$; always strip non-operating items down to core operating earnings.
  \item \textbf{Heavy CapEx hidden by EBITDA.} Compare a B2B software vendor ($\text{EBITDA} =$ \$200M, annual $\text{CapEx} =$ \$10M) against a heavy-equipment rental fleet or regional fiber builder ($\text{EBITDA} =$ \$200M, annual maintenance $\text{CapEx} =$ \$140M). At $\EV =$ \$1,600M, both screen at $8.0\times$ $\EV / \text{EBITDA}$, yet pre-tax cash flow after CapEx is \$190M ($8.4\times$) for the software vendor versus \$60M ($26.7\times$) for the equipment fleet. Treating EBITDA as distributable cash when assets physically wear out every five years confuses gross cash inflow with net wealth creation.
  \item \textbf{Stock-Based Compensation (SBC) excluded from ``Adjusted'' multiples.} Technology companies routinely report ``Non-GAAP EPS'' and ``Adjusted EBITDA'' that add back employee stock grants because no cash leaves the bank account on grant day. Instead, new shares are printed, diluting existing shareholders' fractional ownership of all future cash flows. Valuing a company on $18\times$ Non-GAAP earnings when $35\%$ of those ``earnings'' were handed to employees in equity is equivalent to ignoring a major payroll line item.
\end{enumerate}
\end{trapbox}

In Chapter~\ref{ch:filings-and-profitability}, we take these exact tools to primary financial documents -- learning the universal accounting adjustments for segments, partner revenue share, equity dilution, capital expenditure waves, cohort retention, and private owner add-backs before applying them to public filings, S-1s, and private-company diligence files.

\section*{Exercises}
\addcontentsline{toc}{section}{Exercises}

\begin{exercise}[Fibonacci's contract comparison and continuous compounding]
A supplier offers you two ways to settle an equipment purchase when your opportunity cost of capital is $8.0\%$ per year:
\begin{itemize}
  \item \textbf{Schedule A:} Pay \$420,000 immediately at $t = 0$.
  \item \textbf{Schedule B:} Pay \$150,000 at $t = 0$, \$150,000 at $t = 1$, and \$160,000 at $t = 2$.
\end{itemize}
(a) Compute the present value of Schedule~B under discrete annual compounding at $r = 8.0\%$ and state which schedule is cheaper in present-value terms.
(b) Recompute the present value of Schedule~B under continuous compounding ($d_t = e^{-0.08t}$). Why does continuous compounding make future payments slightly smaller in present value than annual compounding at the same quoted rate $r$?
\end{exercise}

\begin{exercise}[$\IRR$ pathologies versus $\NPV$]
(a) A lithium brine extraction project requires \$200 million at $t = 0$, generates \$620 million of net cash inflow at $t = 1$, and incurs \$462 million of mandatory aquifer restoration at $t = 2$. Solve $\NPV(\IRR) = 0$ algebraically to find both exact $\IRR$ values. Evaluate $\NPV(r)$ at $r = 15\%$ and at $r = 4\%$, and state whether the project should be accepted at each cost of capital.
(b) You can invest in only one of two mutually exclusive cloud-migration contracts at $r = 12\%$: Contract~X costs \$50,000 at $t=0$ and returns \$90,000 at $t=1$; Contract~Y costs \$800,000 at $t=0$ and returns \$1,120,000 at $t=1$. Compute $\IRR$ and $\NPV(0.12)$ for both contracts, plus the incremental $\IRR$ of $(Y - X)$, and identify the wealth-maximizing choice.
\end{exercise}

\begin{exercise}[Two-stage DCF and terminal-value sensitivity]
A vertical SaaS company just generated $\FCF_0 =$ \$25.0 million. You forecast $\FCF$ to grow at $20\%$ per year for $T = 4$ years, then slow to a perpetual steady-state growth rate $g_\infty = 3.0\%$ from year $5$ onward. The discount rate is $r = 9.0\%$.
(a) Compute $\FCF_1, \FCF_2, \FCF_3, \FCF_4$, the present value of the explicit 4-year stage, the year-4 terminal value $V_4 = \frac{\FCF_4(1 + g_\infty)}{r - g_\infty}$, its present value at $t = 0$, and total Enterprise Value $V_0$. What percentage of $V_0$ comes from the terminal value?
(b) Recompute $V_0$ if the discount rate rises by $1.0$ percentage point to $r = 10.0\%$ while $g_\infty$ drops to $2.5\%$. By what percentage does $V_0$ fall?
\end{exercise}

\begin{exercise}[Enterprise Value and leverage-neutral multiples]
Two competing logistics software companies each generate \$80 million of Revenue, \$20 million of EBITDA, \$4 million of Depreciation \& Amortization, and face a $25\%$ corporate tax rate:
\begin{itemize}
  \item \textbf{Company Alpha:} $10.0$ million diluted shares trading at \$24.00/share, \$10 million of Total Debt (at $5\%$ annual interest), \$0 Preferred Stock, and \$50 million of Cash \& Equivalents.
  \item \textbf{Company Beta:} $10.0$ million diluted shares trading at \$14.00/share, \$120 million of Total Debt (at $7.5\%$ annual interest), \$0 Preferred Stock, and \$10 million of Cash \& Equivalents.
\end{itemize}
Compute Market Cap, Enterprise Value ($\EV$), GAAP Net Income, TTM P/E, and $\EV / \text{EBITDA}$ for both Alpha and Beta. Explain why Beta has a lower Share Price and a lower P/E ratio than Alpha despite trading at a higher $\EV / \text{EBITDA}$ multiple.
\end{exercise}

\begin{exercise}[Justified Forward P/E across $\ROIC$ regimes]
Three unlevered companies each expect next-year earnings of $E_1 =$ \$5.00 per share and share the same opportunity cost of capital $r = 10.0\%$:
\begin{itemize}
  \item \textbf{Firm 1:} $\ROIC = 25\%$, plowback rate $b = 32\%$.
  \item \textbf{Firm 2:} $\ROIC = 10\%$, plowback rate $b = 50\%$.
  \item \textbf{Firm 3:} $\ROIC = 6\%$, plowback rate $b = 50\%$.
\end{itemize}
For each firm, compute the perpetual earnings growth rate $g = b \cdot \ROIC$, distributable year-1 cash flow per share $\FCF_1 = E_1(1 - b)$, justified share price $P_0$, and justified Forward P/E ($P_0 / E_1$). Verify why Firm~2 and a zero-growth firm ($b = 0$) share the exact same P/E multiple, whereas Firm~3 trades below $10\times$ despite positive earnings growth.
\end{exercise}
